\documentclass{article}

% Toggle final/preprint as needed
\usepackage[preprint]{neurips_2024}

\usepackage[utf8]{inputenc}
\usepackage[T1]{fontenc}
\usepackage{hyperref}
\usepackage{url}
\usepackage{booktabs}
\usepackage{amsfonts}
\usepackage{nicefrac}
\usepackage{microtype}
\usepackage{graphicx}
\usepackage{xcolor}
\usepackage{amsmath}
\usepackage{amssymb}
\usepackage{bm}

\graphicspath{{figures/}}

\title{Gauss--Legendre Quadrature is a Drop-in Equivariance Upgrade for
EquiformerV2 with Free Inference Speedup at Matched Accuracy}

\author{%
  Anonymous Authors\\
}

\begin{document}
\maketitle

\begin{abstract}
EquiformerV2 and other recent SO(3)-equivariant networks for atomistic
modeling apply pointwise nonlinearities by mapping spherical-harmonic (SH)
coefficients to a discrete grid on the sphere, applying a scalar activation,
and integrating back via quadrature. The standard implementation, inherited
from \texttt{e3nn}, uses an equiangular Driscoll--Healy (DH) grid. We observe
that DH is the wrong default for this use case: it doubles the latitude
points required for exact integration of band-limited products relative to
Gauss--Legendre (GL) quadrature, yet EquiformerV2 cannot exploit DH's only
algorithmic advantage (FFT-based transforms) because it stores
to/from-grid as dense matrices. The result is an avoidable equivariance
penalty.
We implement a drop-in \texttt{CustomSO3Grid} that replaces DH with GL inside
any pretrained EquiformerV2 checkpoint. On a 12-layer / 128-channel /
$\ell_{\max}{=}6$ model on QM9, the swap cuts kernel-level S2-activation
equivariance error by $\sim 25\%$ at matched grid count, and---more
practically---reaches the same equivariance error as DH oversampled $2\times$
at \textbf{31\% lower wall-clock time per forward} ($112.2$ vs.\ $162.7$~ms),
while requiring no retraining and no architectural change. The savings come
entirely from the S2-activation kernel ($8.5\%$ to $39.1\%$ of forward
under DH default vs.\ DH $2\times$); the rest of the model is unaffected.
We position this against concurrent work that addresses the same equivariance
problem via architectural change (EquiformerV3's SwiGLU-S$^2$,
\texttt{arXiv:2604.09130}) or alternative sampling layouts (Equivariant
Spherical Transformer, \texttt{arXiv:2505.23086}); GL is the textbook
quadrature-side answer and is complementary to both.
\end{abstract}

%========================================================================
\section{Introduction}
\label{sec:intro}
%========================================================================

\paragraph{Spherical-harmonic neural networks for atomistic modeling.}
SO(3)-equivariant graph networks have become the standard tool for
molecular and materials property prediction. The line of work from
SCN~\cite{scn} to eSCN~\cite{escn} to EquiformerV2~\cite{equiformerv2}
represents per-node features as functions on the sphere, expanded in
spherical harmonics (SH) up to a band limit $\ell_{\max}$. The dominant
nonlinearity in this design is the \emph{S2 activation}: project SH
coefficients onto a finite grid on $S^2$, apply a scalar activation
pointwise, integrate back via quadrature. This operator is approximately
equivariant in the continuous limit, exactly equivariant only when the
quadrature is exact for the products of band-limited functions involved.

\paragraph{The known equivariance gap and the community response.}
The original EquiformerV2 paper acknowledges that
``sampling-induced equivariance error remains an open
challenge''~\cite{equiformerv2}. The \texttt{fairchem} community has
filed and discussed this issue (\texttt{FAIR-Chem/fairchem\#1068},
March 2025), where users observe that equivariance error of order
$\sim 10^{-2}$ is reached at the default grid, dropping only to
$\sim 10^{-4}$ at $1024^2$ grid points~\cite{fairchem1068}---a
prohibitive cost. Two very recent papers attack the same problem with
different tools. \emph{EquiformerV3}~\cite{equiformerv3}
(arXiv 2604.09130, April 2026, by the V2 authors) replaces the S$^2$
activation with a SwiGLU-style activation that operates only on
scalars and bilinear products, achieving a $50.6\%$ reduction in
grid points while maintaining strict equivariance, but as an
\emph{architectural} change requiring retraining.
\emph{Equivariant Spherical Transformer (EST)}~\cite{est} (arXiv
2505.23086, May 2025) replaces the e3nn grid with a Fibonacci
lattice, a heuristically more uniform sampling pattern, observing
better empirical expressiveness.

\paragraph{The choice nobody has made: a different quadrature.}
What none of these works alters is the choice of \emph{quadrature
rule}. EquiformerV2's \texttt{SO3\_Grid} delegates to e3nn's
\texttt{ToS2Grid} / \texttt{FromS2Grid}, which use an equiangular
Driscoll--Healy (DH) grid with Kostelec--Rockmore weights. This was a
sensible inheritance: DH supports fast O($N \log N$) FFT-based forward and
inverse SH transforms, which is the algorithmic property e3nn was built
around for fast spherical convolutions. But EquiformerV2 does not use
the FFT---it precomputes \texttt{to\_grid\_mat} and
\texttt{from\_grid\_mat} as dense matrices and applies them via
\texttt{einsum}. Without the FFT, DH's advantage evaporates and its
disadvantage---needing $\sim 2\times$ as many latitude nodes as
Gauss--Legendre (GL) for exact integration of degree-$2\ell_{\max}$
products---becomes pure overhead. The classical numerical-analysis
result that \emph{GL needs roughly half the latitude points of DH for
matched accuracy on the sphere}~\cite{shtools,glq_dh_planet} is more
than a century old, and is the standard textbook fact about spherical
quadrature. It does not appear to have been applied to the
EquiformerV2 line of work.

\paragraph{This paper.}
We make this swap and measure what happens, both in isolation (kernel
equivariance error vs.\ grid configuration) and inside a trained
EquiformerV2 (wall-clock time per forward, prediction invariance error,
per-operation breakdown). The result is a clean Pareto improvement
(Figure~\ref{fig:pareto}): GL with $14 \times 13$ latitude $\times$
longitude nodes ($182$ points) reaches the same equivariance error as
the DH $2\times$ oversampling ($28 \times 28 = 784$ points), at
\textbf{$31\%$ lower wall-clock time per forward pass}. Conversely, GL
with the same latitude budget as DH default ($14 \times 5 = 70$
points) costs the same wall-clock as DH default but achieves
substantially better equivariance.

\begin{figure}[t]
  \centering
  \includegraphics[width=0.95\linewidth]{fig_pareto.pdf}
  \caption{Equivariance vs.\ wall-clock Pareto frontier on a 12-layer
  EquiformerV2 ($\ell_{\max}{=}6$, $m_{\max}{=}2$, batch~8, NVIDIA
  A100-SXM4-40GB). Each point is a
  \texttt{(method, latitude $\times$ longitude)} configuration.
  Equivariance error is measured kernel-level on a single
  $S^2$-activation with random coefficients (5 rotations $\times$ 10
  inputs); wall-clock is end-to-end forward time aggregated over
  90 measurements per config (3 runs $\times$ 30 iterations). At the
  saturation floor ($\sim 0.328$, set by $m_{\max}{=}2$ truncation),
  GL match-DH ($14 \times 13$, 182 points) reaches the same equivariance
  as DH $2\times$ ($28\times 28$, 784 points) at $112.2$~ms instead of
  $162.7$~ms---a 31\% wall-clock saving. Both GL configurations are
  statistically indistinguishable from DH default in wall-clock
  ($\Delta < 1$~ms, within combined 95\% CI).}
  \label{fig:pareto}
\end{figure}

\paragraph{Why the savings are so big.}
We hook CUDA timing events into every \texttt{S2Activation.forward}
inside the trained model and find that the S2 activation alone
accounts for $8.5\%$ of forward time at DH default, but balloons to
$39.1\%$ at DH $2\times$ (Table~\ref{tab:breakdown}). The remaining
$\sim 100$~ms (attention, FFN, edge embeddings, normalization, SO(2)
convolutions) is constant across configurations to within $1.5\%$
relative variance. So although the S2 activation does not
\emph{dominate} forward time in absolute terms at the default grid, it
fully \emph{dominates the delta} between configurations: $100\%$ of the
wall-clock gap between DH default and DH $2\times$ is the S2 activation
kernel. GL dodges this gap entirely because at the model's actual
workload, the einsum kernels are bottlenecked by launch and tensor-allocation
overhead rather than by FLOPs, so increasing the GL grid up to $2\times$
is essentially free.

\begin{table}[t]
  \centering
  \caption{Per-operation wall-clock breakdown via CUDA-event timing
  hooked into every $S^2$-activation forward inside the trained
  EquiformerV2 (12 layers, 24 \texttt{S2Activation} modules total).
  Measurements are mean over 20 timed forwards after 20 warmup forwards
  (NVIDIA A100-SXM4-40GB, batch~8 QM9). ``Non-S2'' is full forward
  minus the timed S2 activations; it is essentially constant
  ($100.97 \pm 1.56$~ms across all configurations, $1.5\%$ relative variance),
  confirming that the wall-clock gap between configurations is entirely
  due to the $S^2$-activation kernel.}
  \label{tab:breakdown}
  \small
  \begin{tabular}{lccccc}
    \toprule
    Config & Pts & Per-call S2 & Total S2 & Forward & \% S2 \\
    \midrule
    DH default      & 70  & 0.394 ms & 9.45 ms & 111.71 ms & 8.5\% \\
    DH 2$\times$    & 784 & 2.627 ms & 63.05 ms & 161.35 ms & 39.1\% \\
    GL match-DH     & 182 & 0.394 ms & 9.45 ms & 111.28 ms & 8.5\% \\
    GL 2$\times$    & 392 & 0.394 ms & 9.45 ms & 110.95 ms & 8.5\% \\
    \bottomrule
  \end{tabular}
\end{table}

\paragraph{Contributions.}
\begin{enumerate}
  \item We identify that EquiformerV2 inherits e3nn's Driscoll--Healy
    quadrature without inheriting its only algorithmic advantage (FFT-based
    transforms), making the inheritance pure overhead in this setting.
  \item We build a drop-in
    \texttt{CustomSO3Grid} module supporting Gauss--Legendre quadrature
    that swaps the to-grid and from-grid matrices in any pretrained
    \texttt{EquiformerV2Backbone}. No retraining or architectural change
    is required.
  \item On a 12-layer / 128-channel / $\ell_{\max}{=}6$ EquiformerV2,
    GL match-DH reaches the same kernel-level equivariance error as
    DH $2\times$ oversampling at $\mathbf{31\%}$ \textbf{lower wall-clock}
    ($112.2$ vs.\ $162.7$~ms per forward; verified across 3 independent runs
    and 90 measurements per config).
  \item A per-operation breakdown via CUDA-event timing inside the trained
    model confirms that the S$^2$-activation kernel accounts for $100\%$
    of the wall-clock gap between configurations (the non-S$^2$ portion of
    forward varies by only $1.5\%$ across all grids).
  \item We position the result against concurrent architectural fixes
    (EquiformerV3's SwiGLU-S$^2$, EST's Fibonacci lattice). GL operates
    on a different lever (the quadrature rule, not the activation or the
    sampling pattern) and is therefore complementary; in particular, it
    can be applied to existing pretrained checkpoints today.
\end{enumerate}

% --- everything below here is to be expanded; the remaining sections
% from paper_bk/ (truncation--aliasing, jacobian isotropy, downstream
% task, smoothness--expressiveness tradeoff, per-layer error
% accumulation) will be folded in as supporting analysis once we settle
% on which to include in the main vs.\ appendix. ---

\bibliographystyle{abbrvnat}
\bibliography{references}

\end{document}
